\section{GIỚI THIỆU BÀI TOÁN}
\subsection{Đặt vấn đề}
Trong các hệ thống hỗ trợ lái xe và giao thông thông minh, biển báo cung cấp những thông tin trực tiếp về giới hạn tốc độ, hướng di chuyển và các tình huống cần cảnh báo. Khả năng nhận dạng đúng biển báo phụ thuộc trước hết vào chất lượng hình ảnh mà hệ thống tiếp nhận. Khi điều kiện chiếu sáng suy giảm, những đặc trưng quan trọng như đường viền, chữ số và mũi tên có thể trở nên khó phân biệt, ngay cả khi vị trí của biển báo đã được xác định.

Một hướng giải quyết thường được cân nhắc là tăng cường ảnh trước khi đưa vào mô hình nhận dạng. Tuy nhiên, tăng độ sáng không đồng nghĩa với phục hồi đầy đủ thông tin. Một phép biến đổi có thể làm vùng tối dễ quan sát hơn nhưng đồng thời khuếch đại nhiễu, thay đổi màu sắc hoặc làm nhòe những chi tiết nhỏ. Đối với biển báo giao thông, chỉ một nét bị mất trong chữ số hoặc một hướng mũi tên bị biến dạng cũng có thể dẫn đến dự đoán sai lớp.

Nhóm đặt ra câu hỏi chính cho đề tài là: \textit{các phương pháp xử lý ảnh có giúp mô hình nhận dạng biển báo chính xác hơn khi ảnh bị thiếu sáng hay không?} Để trả lời câu hỏi này, nhóm không chỉ quan sát ảnh trước và sau xử lý mà còn tính các chỉ số chất lượng ảnh và kết quả phân loại. Những phép so sánh này cần được thực hiện trên cùng dữ liệu và cùng mô hình nhận dạng.

Khi tìm hiểu bài toán, nhóm nhận thấy ảnh thiếu sáng cần được xem xét cùng với nhiễu. Vùng tối có tín hiệu nhỏ nên nhiễu có thể chiếm tỷ lệ đáng kể so với phần thông tin hữu ích. Nếu phương pháp chỉ nâng độ sáng, nó có thể làm rõ cả cấu trúc biển báo lẫn nhiễu. Vì vậy, nhóm cho rằng việc chỉ làm tối ảnh bằng Gamma là chưa đủ để khảo sát tình huống ảnh vừa tối vừa nhiễu.

Nhóm cũng phân biệt chất lượng ảnh với khả năng nhận dạng của mô hình. PSNR phản ánh sai số cường độ pixel, còn SSIM xem xét sự tương đồng về cấu trúc cục bộ. Trong khi đó, mô hình phân loại dựa vào các đặc trưng đã học để phân biệt từng lớp. Một ảnh có sai số trung bình thấp vẫn có thể bị mất nét chữ số. Ngược lại, một ảnh chưa gần ảnh tham chiếu về độ sáng vẫn có thể giữ được hình dạng cần cho nhận dạng.

Để so sánh các phép xử lý, nhóm cần giữ các yếu tố còn lại ổn định. Nếu nhóm vừa đổi cách xử lý ảnh vừa huấn luyện lại mô hình riêng cho mỗi phương pháp, kết quả sẽ phản ánh đồng thời cả hai thay đổi. Vì vậy, nhóm sử dụng một kiến trúc ResNet-18, huấn luyện trên ảnh gốc rồi giữ nguyên trọng số khi đánh giá. Cách làm này cho biết mỗi phép xử lý giúp ích như thế nào đối với bộ nhận dạng đã có, nhưng chưa cho biết kết quả khi mô hình được huấn luyện thích nghi với ảnh sau xử lý.

\subsection{Yêu cầu và phạm vi nghiên cứu}
Đề bài yêu cầu khảo sát ảnh hưởng của xử lý ảnh đến hiệu năng nhận dạng biển báo giao thông trong điều kiện thiếu sáng trên bộ dữ liệu GTSRB. Ảnh suy giảm được mô phỏng bằng Gamma kết hợp nhiễu Gaussian và Poisson ở nhiều mức độ. Các phương pháp tăng cường cần so sánh gồm hiệu chỉnh Gamma, cân bằng histogram (HE), cân bằng histogram cục bộ có giới hạn tương phản (CLAHE), Retinex, Zero-DCE, Retinexformer và phương pháp dựa trên diffusion.

Ngoài các phương pháp tăng cường ảnh, nhóm khảo sát ba bộ lọc Gaussian, bilateral và Non-local Means. Nhóm thử từng bộ lọc riêng lẻ, sau đó kết hợp chúng với Gamma và CLAHE theo hai thứ tự. Độ mạnh của bộ lọc được chọn trên tập validation, tức tập dùng để lựa chọn cấu hình, và giữ giống nhau khi đảo thứ tự. Nhóm dùng PSNR và SSIM để đánh giá chất lượng ảnh, đồng thời dùng Accuracy và F1-score để đánh giá nhận dạng. Macro-F1 là tiêu chí chính vì chỉ số này cho mỗi lớp vai trò như nhau. Từ những kết quả đó, nhóm phân tích mối liên hệ giữa chất lượng ảnh và khả năng nhận dạng để lựa chọn quy trình tiền xử lý.

Trong phạm vi bài tập lớn, nhóm chỉ xử lý vùng ảnh đã chứa biển báo, không xử lý toàn cảnh giao thông. Bộ phân loại đưa mỗi ảnh vào một trong 43 lớp của GTSRB. Hệ thống chưa có bước phát hiện vị trí biển báo, theo dõi đối tượng hoặc xử lý video. Nhóm tạo điều kiện thiếu sáng từ ảnh gốc bằng mô hình suy giảm đã xác định. Nhờ vậy, mỗi ảnh suy giảm có một ảnh tham chiếu tương ứng để tính PSNR và SSIM.

\subsection{Đầu vào và đầu ra của hệ thống}
\textbf{Đầu vào:} Hệ thống nhận ảnh màu RGB đã cắt vùng quan tâm và đổi về kích thước $64\times64$ pixel. Giá trị cường độ được chuẩn hóa về $[0,1]$. Từ ảnh gốc $X$, nhóm tạo ảnh thiếu sáng $Y$ theo mức suy giảm và hạt giống sinh nhiễu đã chọn.

\textbf{Đầu ra của khối xử lý ảnh:} Mỗi phương pháp tạo ảnh $\widehat{X}=E(Y)$ có cùng kích thước, thứ tự kênh màu và miền giá trị với đầu vào. Hàm $E$ chỉ nhận ảnh thiếu sáng. Nhóm không đưa ảnh tham chiếu hoặc nhãn lớp vào thuật toán tăng cường.

\textbf{Đầu ra của khối nhận dạng:} ResNet-18 trả về vector 43 giá trị dự đoán, từ đó nhóm lấy lớp có giá trị lớn nhất làm kết quả phân loại. Trước khi vào mạng, ảnh được nội suy lên $224\times224$ và chuẩn hóa theo ImageNet. Bước này phục vụ riêng bộ nhận dạng. Các phép xử lý ảnh và phép đo PSNR, SSIM vẫn được thực hiện ở $64\times64$.

\textbf{Đầu ra nghiên cứu:} Báo cáo trình bày các bảng PSNR, SSIM, Accuracy, Macro-F1 và thời gian xử lý. Nhóm bổ sung hình ảnh trước và sau xử lý, phân tích từng thành phần suy giảm bằng ablation và khảo sát tương quan giữa các chỉ số. Quy trình đề xuất được chọn trên validation trước khi đánh giá test.

\subsection{Các hướng tiếp cận và giải pháp khảo sát}
Các phương pháp truyền thống xử lý trực tiếp cường độ hoặc phân bố mức xám của ảnh. Gamma biến đổi từng pixel bằng hàm lũy thừa. HE phân bố lại mức sáng theo histogram toàn ảnh, còn CLAHE thực hiện theo từng vùng nhỏ. Retinex dựa trên mối quan hệ giữa phản xạ và chiếu sáng để làm rõ nội dung. Những phương pháp này có cách triển khai tương đối dễ theo dõi và thường có chi phí tính toán thấp.

Nhóm phương pháp học sâu sử dụng thông tin đã học để điều chỉnh hoặc phục hồi ảnh. Zero-DCE ước lượng đường cong tăng sáng, Retinexformer kết hợp biểu diễn chiếu sáng với Transformer, còn phương pháp diffusion được khảo sát thực hiện phục hồi có điều kiện trong miền wavelet. Dự án sử dụng trọng số công bố của tác giả, không huấn luyện lại các bộ tăng cường trên GTSRB~\cite{zerodce}, \cite{retinexformer}, \cite{diffusion}.

Giải pháp của bài tập lớn không giả định một nhóm phương pháp luôn tốt hơn nhóm còn lại. Tất cả được đánh giá trên cùng ảnh suy giảm, cùng tập mẫu và cùng các checkpoint ResNet-18, tức các bộ trọng số đã lưu sau huấn luyện. Mốc đối chứng là \term{Identity}, tức ảnh tối giữ nguyên. Ảnh gốc \term{Clean} được dùng làm tham chiếu và kiểm tra hiệu năng nhận dạng khi chưa áp dụng mô hình suy giảm.

\subsection{Đóng góp triển khai và xây dựng quy trình thực nghiệm}
\label{sec:contributions}
Nhóm xây dựng quy trình chuẩn bị dữ liệu từ các archive GTSRB chính thức, gồm đọc ảnh PPM và annotations, chuyển ảnh sang RGB, cắt vùng biển báo theo ROI với padding 5\% và nội suy song tuyến tính về kích thước $64\times64$. Các bước này vận dụng kiến thức về biểu diễn ảnh, hệ tọa độ, phép cắt và lấy mẫu lại. Quy trình này đồng thời tạo đầu vào thống nhất cho tất cả phương pháp tăng cường và bộ phân loại. Chi tiết cài đặt và ý nghĩa của các lựa chọn được trình bày trong phần xây dựng tập dữ liệu.

Nhóm dùng ảnh số thực trong miền $[0,1]$ khi tạo suy giảm và tăng cường. Trong khảo sát tăng sáng, kết quả trung gian được lưu bằng NumPy. Với khử nhiễu và kết hợp, mã truyền trực tiếp mảng float32 trong bộ nhớ để hạn chế lượng tử hóa và giảm thao tác đọc ghi. Riêng HE, CLAHE và NLM có bước chuyển sang ảnh 8 bit theo cài đặt OpenCV. Nhóm nêu rõ các khác biệt này ở phần phương pháp vì chúng có thể ảnh hưởng đến so sánh chất lượng ảnh và phân loại.

Nhóm kết hợp quy trình xử lý ảnh với kiểm tra trùng lặp, chia dữ liệu theo track và lưu checksum cùng thông tin nguồn để hỗ trợ tái lập thí nghiệm. Giá trị của phần triển khai nằm ở khả năng kiểm soát, giải thích và kiểm chứng toàn bộ luồng dữ liệu. Các thao tác tải, giải nén và gom annotations là các bước hỗ trợ. Đây là đóng góp về triển khai và thiết kế thực nghiệm, không phải đề xuất thuật toán xử lý ảnh mới. Padding, kích thước crop và phép nội suy được trình bày là các quyết định thiết kế. Nhóm chưa có thí nghiệm so sánh riêng để khẳng định chúng là cấu hình tối ưu.

\subsection{Ý nghĩa và bố cục báo cáo}
Về kiến thức xử lý ảnh, đề tài kết nối phép biến đổi điểm, histogram, lọc Gaussian, mô hình nhiễu và các chỉ số chất lượng trong một quy trình thống nhất. Về thực nghiệm, đề tài làm rõ sự đánh đổi giữa độ sáng, mức nhiễu, tính toàn vẹn của chi tiết và chi phí suy luận. ResNet-18 đóng vai trò công cụ kiểm chứng thông tin được giữ lại sau xử lý, thay vì là đối tượng tối ưu chính của báo cáo.

Báo cáo bắt đầu bằng cách chuẩn bị dữ liệu và tạo ảnh suy giảm. Phần tiếp theo giải thích cơ sở và cách cài đặt các phương pháp xử lý ảnh. Sau đó, nhóm trình bày thiết lập thí nghiệm, kết quả nhận dạng và chất lượng ảnh. Phần thảo luận tập trung vào tương quan, ảnh hưởng của nhiễu và những trường hợp nhận dạng sai. Cuối cùng, nhóm đề xuất quy trình tiền xử lý và nêu các kết luận cùng giới hạn của đề tài.
